\documentclass[a4paper,12pt]{article}
\usepackage{graphicx}
\usepackage[utf8]{inputenc}  
\usepackage[T1]{fontenc}
\usepackage{gensymb}
\usepackage{lscape}

\usepackage{helvet}%Police Arial
\renewcommand{\familydefault}{\sfdefault}

\usepackage[top = 1cm, bottom = 2cm, left = 1cm, right = 1cm]{geometry}
\usepackage{titlesec}
\setcounter{secnumdepth}{4}
\usepackage{xcolor}
\usepackage{amssymb} %double flèche de réaction
\usepackage{xtab}%tableau sur plusieurs pages
\usepackage{esvect}%grands vecteurs

\usepackage[version=4]{mhchem} %equations chimiques

\usepackage{array}
\usepackage{chemfig}
\usepackage{multirow}
\usepackage{sectsty}
\usepackage{caption}
\usepackage{amsmath}
\usepackage{amssymb}
\usepackage[cal=boondoxo]{mathalfa} 
\usepackage{enumitem}
\usepackage{lastpage}
\usepackage{tcolorbox}
\usepackage{siunitx}
\usepackage[european, straightvoltages]{circuitikz}
\usepackage{tikz}
\usetikzlibrary{shapes.geometric}
\usetikzlibrary{calc}
\usetikzlibrary{automata,positioning}
\usepackage{multicol}
\usepackage{eqnarray}
\usepackage{tabularx,colortbl}%colorer fond cellules

\sisetup{
	detect-all,
	output-decimal-marker={,},
	group-minimum-digits = 3,
	group-separator={~},
	number-unit-separator={~},
	inter-unit-product={~}
} %setup pour SIunitx

\usepackage{listings}%pour insérer le code python
\definecolor{codegreen}{rgb}{0,0.6,0}
\definecolor{codegray}{rgb}{0.5,0.5,0.5}
\definecolor{codepurple}{rgb}{0.58,0,0.82}
\definecolor{backcolour}{rgb}{0.95,0.95,0.92}

\graphicspath{ {/home/remy/Documents/Enseignement/2022/Tspe/C11/Ressources/} }

\sectionfont{\fontsize{14}{8}\selectfont}
\subsectionfont{\fontsize{14}{8}\selectfont}
\renewcommand{\thesection}{\arabic{section}.}
\renewcommand{\thesubsection}{\hspace{1cm}\alph{subsection}.}

\title{\vspace{-1cm}\large Chapitre 10 : Systèmes électriques capacitifs}
\date{}

\newpagestyle{mystyle}{%
	\setfoot{}{\thepage\,$|$\,\pageref{LastPage}}{}
}
\renewpagestyle{plain}{%
	\setfoot{}{\thepage\,$|$\,\pageref{LastPage}}{}
}

\pagestyle{mystyle}

\newenvironment{itemize*}%
{\begin{itemize}[label=$-$]%
		\setlength{\itemsep}{0pt}%
		\setlength{\parskip}{0pt}}%
	{\end{itemize}}

\newenvironment{itemize**}%
{\begin{itemize}[label=]%
		\setlength{\itemsep}{0pt}%
		\setlength{\parskip}{0pt}}%
	{\end{itemize}}

\begin{document}
	\maketitle
	\vspace{-2.5cm}


\section*{Ce qu'il faut connaitre :}
\begin{multicols}{2}
	\subsection*{Le vocabulaire}
	\begin{itemize}
		\item[$\square$] comportement capacitif
		\item[$\square$] condensateur
	\end{itemize}
	\columnbreak
	\subsection*{Les grandeurs et relations associées}
	\begin{itemize}
		\item[$\square$] intensité
		\item[$\square$] tension
		\item[$\square$] charge
		\item[$\square$] capacité
	\end{itemize}
\end{multicols}
	

\renewcommand{\arraystretch}{2}

\section*{Ce qu'il faut savoir faire :}
\begin{table}[h!]
	\begin{tabular}{|m{12cm}|c|c|}
		\hline \rowcolor{gray!20}\centering \textbf{La compétence }& \textbf{Où la trouver}& \textbf{Je la maitrise ?}\\
		\hline Relier l'intensité d'un courant électrique au débit de charges. & &\\
		\hline Identifier des situations variées où il y a accumulation de charges de signes opposés sur des surfaces en regard. & & \\
		\hline Citer des ordres de grandeur de valeurs de capacités usuelles. & &\\
		\hline Identifier et tester le comportement capacitif d'un dipôle. &&\\
		\hline Illustrer qualitativement, par exemple à l'aide d'un microcontrôleur, d'un multimètre ou d'une carte d'acquisition, l'effet de la géométrie d'un condensateur sur la valeur de sa capacité.&&\\
		\hline Établir et résoudre l'équation différentielle vérifiée par la tension aux bornes d'un condensateur (charge et décharge).&&\\
		\hline Expliquer le principe de fonctionnement de quelques capteurs capacitifs.&&\\
		\hline Étudier la réponse d'un dispositif modélisé par un dipôle RC.&&\\
		\hline Déterminer le temps caractéristique d'un dipôle RC à l'aide d'un microcontrôleur, d'une carte d'acquisition ou d'un oscilloscope.&&\\
		\hline \underline{Capacité mathématique :} Résoudre une équation différentielle linéaire du premier ordre à coefficient constant avec un second membre constant.&&\\
		\hline
	\end{tabular}
\end{table}
\end{document}